\section{The third factorial cumulant and rational triples}
\label{p13:third:section}

The quadratic reference removes the first two factorial cumulants.  To
identify the remaining third-order contribution, two tasks remain distinct:
we must bound every compatible triple, and we must show that the positive
rational examples survive in the signed cumulant.  The incompatible
activity estimate in Section~\ref{p13:activity:section} supplies the second
step.  The first step combines rational overlap geometry with the
nonaligned pair bound proved above.

Recall the exact marginals $p_x=\E J_x$ and the relative nullity
\[
 \rho_{\mathrm{rel}}(T)=\sum_{x\in T}\dim S_x
                    -\dim\Bigl(\sum_{x\in T}S_x\Bigr).
\]
Write $d_{xy}=\dim(S_x\cap S_y)$ and
$c(T)=\E\prod_{x\in T}(J_x-p_x)$.  A triple is connected when its
spaces have no nontrivial direct-sum partition.  For compatible triples,
$q(T)=2^{\rho_{\mathrm{rel}}(T)}\prod_{x\in T}p_x$ and
$0\leq c(T)\leq q(T)$, by the identities established in
Section~\ref{p13:activity:section}.

The third factorial cumulant and its optimal exponent are stated in
Theorem~\ref{p13:third:main} of the article. We give the geometric
and probabilistic proof here.

\subsection{The largest canonical overlap of three rational windows}

Consider integer slopes $u<v<w$ with $(u,v,w)=1$ and $w<2u$.
In the common parameter, the three windows are intervals of lengths
$L/u,L/v,L/w$ on the lattices $u^{-1}\Z$, $v^{-1}\Z$, and
$w^{-1}\Z$.  Translation of a window moves its interval and leaves
the lattice origin fixed.  The densities of pairwise coincidence
lattices are $(u,v),(u,w),(v,w)$; the density of the triple coincidence
lattice is one.

Let $I_u,I_v,I_w$ denote the three parameter intervals.  Up to a bounded
endpoint error, the number of dimensions lost when equal values are
merged is at most
\begin{align*}
 &(u,v)|I_u\cap I_v|+(u,w)|I_u\cap I_w|
       +(v,w)|I_v\cap I_w|-|I_u\cap I_v\cap I_w|\\
 &\quad=(u,v)|I_u\cap I_v|+((u,w)-1)|I_u\cap I_w|\\
 &\qquad\quad+((v,w)-1)|I_v\cap I_w|
                         +|I_w\cap(I_u\cup I_v)|.
\end{align*}
All coefficients in the last expression are nonnegative.  Bounding its
four lengths by $L/v,L/w,L/w,L/w$ gives
\begin{equation}
 D(u,v,w)L+O(1),\qquad
 D(u,v,w)=\frac{(u,v)}v+\frac{(u,w)+(v,w)-1}{w}.
 \label{p13:third:density}
\end{equation}
Nested intervals attain the coefficient $D$.
The endpoint error is absolute: each of the four lattice counts has a
bounded rounding error, even when the integer slopes vary.

\begin{lemma}[A sharp rational overlap inequality]
\label{p13:third:rational-rank}
For all integers $u<v<w<2u$ with $(u,v,w)=1$,
\begin{equation}
 D(u,v,w)\leq\frac{19}{30}.
 \label{p13:third:1930}
\end{equation}
Equality holds for $(u,v,w)=(10,12,15)$.
This is an inequality for the canonical coincidence geometry; arithmetic
relations among the distinct normalized values require a separate bound.
\end{lemma}

\begin{proof}
If the reduced height of every pair is at least five, each pairwise
intersection density is at most $1/5$.  Dropping the subtracted triple
intersection in \eqref{p13:third:density} gives $D\leq3/5$.
Otherwise a pair has reduced ratio $2:3$ or $3:4$.  Reduced heights one
and two are impossible for distinct slopes whose extreme ratio is
strictly less than two.

For a pair $2:3$, write the three slopes before reordering as
$(2d,3d,c)$, with $(c,d)=1$ and $3/2<c/d<4$.
The pair itself contributes at most $L/3+O(1)$ lost dimensions.
In the third window, positions that can coincide with either old family
have density at most
\[
 \frac{(c,2)+(c,3)-1}{c}.
\]
This is at most $1/4$ except for the possible numerators $c=9,12$,
after removing physical repetitions and endpoint ratios.
For completeness, if $5\leq c\leq15$, direct evaluation of
$(c,2)+(c,3)-1$ leaves just $9,12$ above $c/4$; the numerator is at most
four for $c\geq16$.  The smaller numerators $2,3$ give physical
repetitions, and $4,6$ have no coprime denominator in the strict
interior range apart from endpoint cases.  The exceptional fractions
and their values of \eqref{p13:third:density} are
\[
 \begin{array}{c|c|c}
 c/d & (u,v,w)&D(u,v,w)\\ \hline
 9/4&(8,9,12)&11/18\\
 9/5&(9,10,15)&17/30\\
 12/5&(10,12,15)&19/30\\
 12/7&(12,14,21)&4/7.
 \end{array}
\]
Every other case has $D\leq1/3+1/4=7/12$.

For a pair $3:4$, write $(3d,4d,c)$ with $(c,d)=1$ and
$2<c/d<6$.  The new density is at most
$((c,3)+(c,4)-1)/c$.  It is at most $1/3$ except for $c=8,12$,
after the same exclusions: evaluating the integers $c\leq17$ suffices,
and the numerator is at most six for $c\geq18$.
The exceptional reduced fractions are $8/3$ and $12/5$, giving
$(8,9,12)$ and $(12,15,20)$, of densities $11/18$ and $3/5$.
All other cases have $D\leq1/4+1/3=7/12$.
The largest value in these cases is $19/30$.
\end{proof}

\subsection{All triples above a strong aligned pair}

A fixed rational pair can be compressed into its true normalized values
even when some of these values lack private prime coordinates.
Injectivity is not needed for the inclusion of event spaces.  This is
what allows the next estimate to cover all pairs in the chosen channels,
including the arithmetically exceptional ones.

\begin{proposition}[Triples meeting finitely many aligned channels]
\label{p13:third:aligned-family}
Fix a finite set $\mathcal F$ of reduced interior channels $a:b$, with
$b\geq3$ and $b/2<a<b$, before $N$.  Then
\begin{equation}
 \sum_{\substack{|T|=3,\ T\ {\rm connected},\ q(T)>0\\
            \text{some pair of }T\text{ is aligned in }\mathcal F}}q(T)
       \ll_{\mathcal F,C,\varepsilon}N^{-41/30+\varepsilon}.
 \label{p13:third:aligned-family-bound}
\end{equation}
The sum allows irregular admissible windows, bad normalization of the
source pair, and contacts at distance exactly $L$.
\end{proposition}

The complete normalization argument is in
Section~\ref{p14:comp:third:aligned-family}.  After exact merging, the
residual kernel has dimension $o(L)$ uniformly even when the last
rational coefficients have polynomial height in $B$.  Distinct
physical slopes then lose at most $19L/30+O(1)$ dimensions through
canonical coincidences, so each triple has probability at most
$N^{-71/30+o(1)}$.  Their $NB^{O(1)}$ hosts give
$N^{-41/30+\varepsilon}$.  The proof separately controls nonaligned
targets, repeated physical slopes, and dyadic endpoint populations;
none of these cases requires good normalization of the source pair.

\begin{corollary}[Strong aligned pairs]
\label{p13:third:aligned}
For every fixed $\theta>0$ and every $\varepsilon>0$,
\begin{equation}
 \sum_{\substack{|T|=3,\ T\ {\rm connected},\ q(T)>0\\
        \exists\{x,y\}\subset T\ {\rm aligned}:\ d_{xy}\geq\theta B}}
             q(T)\ll_{C,\theta,\varepsilon}N^{-41/30+\varepsilon}.
 \label{p13:third:strong-aligned-bound}
\end{equation}
\end{corollary}

\begin{proof}
The uniform bound $d_{xy}\leq L/b+o(L)$ forces $b\leq2/\theta$
for sufficiently large $N$.  Thus only a finite collection of heights
is relevant.  Heights at least three satisfy
Proposition~\ref{p13:third:aligned-family}.
Height two forces the source pair to the dyadic endpoints: there are
$B^{O(1)}$ pairs of weight at most $N^{-3/2+o(1)}$.
Proposition~\ref{p1:fixed:propagation} adds a third window with a loss
$N^{\varepsilon'}$, which is harmless.
At height one, compatible distinct windows have disjoint interiors
and at most one exact unit, so their relative intersection is $o(L)$.
They cannot reach the threshold $\theta B$.
\end{proof}

\subsection{A positive family that survives in the total cumulant}

The equality case $(10,12,15)$ of the overlap inequality also supplies
a population of actual compatible run starts.  Its role is stronger
than merely maximizing a geometric formula: its centered activity is
positive, and there are linearly many admissible translations.

\begin{proposition}[The rational lower bound]
\label{p13:third:lower}
For all sufficiently large $N$ in the critical band,
\begin{equation}
 K_3^+:=\sum_{\substack{|T|=3,\ T\ {\rm connected}\\q(T)>0}}c(T)
                  \gg_C N^{-41/30}.
 \label{p13:third:K3-lower}
\end{equation}
In particular $\kappa_{3,N}\geq c_C N^{-41/30}$ for a suitable
$c_C>0$.
\end{proposition}

\begin{proof}
Take integers $t$ with $N\leq10t<12t<15t<2N$.
There are $N/30+O(1)$ choices.  Multiplication of the three windows by
$6,5,4$ respectively gives the true normalized sets
\[
 U_1=60t+6\{-1,0,\ldots,L-1\},\quad
 U_2=60t+5\{-1,0,\ldots,L-1\},\quad
 U_3=60t+4\{-1,0,\ldots,L-1\}.
\]
Choose $H=10B$, larger than the diameter of their union.
By \eqref{p1:arith:onedfect}, only $N^{1/2+o(1)}$ integers in the
relevant range have $K_H(n)=1$.  For each of the $O(B)$ possible
normalized offsets there is at most one parameter $t$ per such integer.
Hence, outside $N^{1/2+o(1)}$ choices of $t$, every distinct value in
$U_1\cup U_2\cup U_3$ has a private prime above $H$.
Its square-class map is then injective.

The three sets meet at $60t$.  The sum of their even-selection spaces
is consequently the even-selection space of their union.
The pairwise intersection sizes are
\[
 1+\left\lfloor\frac{L-1}{6}\right\rfloor,\qquad
 1+\left\lfloor\frac{L-1}{3}\right\rfloor,\qquad
 1+\left\lfloor\frac{L-1}{5}\right\rfloor,
\]
and the triple intersection has size
$1+\lfloor(L-1)/15\rfloor$.
Inclusion--exclusion therefore gives the exact relative nullity
\begin{equation}
 \rho_L=\left\lfloor\frac{L-1}{6}\right\rfloor
        +\left\lfloor\frac{L-1}{3}\right\rfloor
        +\left\lfloor\frac{L-1}{5}\right\rfloor
        -\left\lfloor\frac{L-1}{15}\right\rfloor
       =\frac{19}{30}L+O(1).
 \label{p13:third:witness-nullity}
\end{equation}

The normalized predecessors have offsets $-6,-5,-4$, each belonging
to just one set.  Assign sign $-1$ to a private prime of each of these
three values and sign $+1$ to every other prime occurring in the
windows and multipliers.  The chosen primes are distinct and occur in
no other normalized value.  Each original predecessor then has sign
$-1$ and every internal value has sign $+1$.  This verifies affine
compatibility explicitly.  For all these parameters,
\begin{equation}
 q(\{10t,12t,15t\})=p^3 2^{\rho_L}.
 \label{p13:third:witness-probability}
\end{equation}

The three pair nullities are the first three floor terms in
\eqref{p13:third:witness-nullity}.  Since all marginals are $p$ on
this private-prime population, the centered identity is exact:
\begin{equation}
 \begin{split}
 c(\{10t,12t,15t\})=p^3\biggl\{2^{\rho_L}
  &-2^{\lfloor(L-1)/6\rfloor}-2^{\lfloor(L-1)/3\rfloor}\\
  &-2^{\lfloor(L-1)/5\rfloor}+2\biggr\}.
 \end{split}
 \label{p13:third:witness-centered}
\end{equation}
The largest subtracted exponent is $L/3+O(1)$, whereas the first
exponent is $19L/30+O(1)$.  Thus $c(T)\geq q(T)/2>0$ for large $N$.
This also ensures connectedness: a direct singleton split would make
the centered coefficient vanish.
All compatible connected triple coefficients are nonnegative, and
the triples obtained from distinct $t$ are different.  Consequently
\[
 K_3^+\gg Np^3 2^{19L/30}
                 \gg_C N^{-41/30}.
\]
Finally \eqref{p13:activity:kappa-comparison} gives
$\kappa_{3,N}=6K_3^++O(N^{-3/2+\varepsilon'})$.
Choose $\varepsilon'<3/2-41/30=2/15$ to make the error negligible.
\end{proof}

When $N=2^L$ and $L=30m$, the contribution of the private-prime
population in this proof to $K_3^+$ is
$(1/120+o(1))N^{-41/30}$, since $\rho_L=19m-2$.
This coefficient belongs to the displayed subfamily; it is not a
coefficient for the full cumulant.

\subsection{Partitioning every compatible triple}

\begin{proof}[Proof of Theorem~\ref{p13:third:main}]
Proposition~\ref{p13:third:lower} supplies the lower bound.
For the upper bound, fix $\eta=1/30$ and then a positive
$\theta<\eta/200$ for which
Proposition~\ref{p13:conf:near-primitive} applies.
All these parameters are fixed before $N$.
We bound $K_3^+$ by partitioning its compatible connected triples.

If $\rho_{\mathrm{rel}}(T)<\eta B$, the host bound of
Corollary~\ref{p13:half:connected} and $p_x\leq2p$ give
\[
 \sum q(T)\ll
        N^{3/2+\varepsilon'}N^{-3+\eta+o(1)}
       =N^{-22/15+\varepsilon'+o(1)}.
\]
If $\rho_{\mathrm{rel}}(T)\geq\eta B$ and every proper pair has
$d_{xy}<\theta B$, Proposition~\ref{p13:conf:near-primitive}
gives total homogeneous mass $O(N^{-2+\theta+\varepsilon'})$,
which bounds the compatible mass as well.

It remains to consider a pair with $d_{xy}\geq\theta B$.
If any such pair is aligned, Corollary~\ref{p13:third:aligned}
gives $O(N^{-41/30+\varepsilon'})$.
Otherwise choose a strong nonaligned pair.  Theorem~\ref{p13:rig:nonaligned}
gives
\[
 d_{xy}\leq\left(\frac23\log2+o(1)\right)L.
\]
For the third space,
\[
 \rho_{\mathrm{rel}}(T)
   =d_{xy}+\dim\bigl(S_z\cap(S_x+S_y)\bigr)
   \leq d_{xy}+L
   \leq\left(1+\frac23\log2+o(1)\right)L.
\]
Since $(2/3)\log2<1/2$, this is eventually at most $3L/2$.
Proposition~\ref{p13:conf:strong-profile}, with $\sigma=3/2$,
therefore bounds this sector by $O(N^{-3/2+\varepsilon'})$.
This is where the unconditional nonaligned-pair rigidity enters the
third-cumulant theorem.

Each exponent other than $41/30$ is strictly larger than $41/30$.
Shrink the internal error margins to obtain
$K_3^+\ll N^{-41/30+\varepsilon}$.
The identity \eqref{p13:activity:kappa-comparison} transfers the estimate to
$\kappa_{3,N}$, completing the upper bound on the full critical band.
\end{proof}

The gap $41/30-4/3=1/30$ explains why third-order activity is below
the optimal quadratic remainder scale.  This conclusion does not
identify the optimal error for a reference that also retains cubic or
quartic terms; such a reference requires its own reconstruction and
lower-bound argument.

The rank threshold $\eta=1/30$ and the weak-pair threshold
$\theta<\eta/200$ are distinct from the accuracy parameters.
Using the convention in Section~\ref{p1:intro:dependencies}, prescribe
$\varepsilon_{\mathrm{out}}>0$ and choose
$0<\tau<\min(\varepsilon_{\mathrm{out}}/4,1/960)$.
Allocate smaller internal losses so that each completed sector loses
at most $\tau$. The nonaligned target above a rational pair retains
$1/120-\tau>0$, and the incompatible-triple error retains
$2/15-\tau>0$. These budgets include offset choices and residual-rank
losses before $N$ is taken large. The latter error is therefore
negligible in the positive rational lower bound as well.

\subsection{All triples above finitely many rational pair channels}
\label{p14:comp:third}

The canonical overlap inequality and the positive $10:12:15$ witness
are proved in Section~\ref{p13:third:section}.  The remaining technical
input is uniformity of the upper bound across every triple above a
fixed finite set of aligned channels, including arithmetically
exceptional source pairs.  We prove that input here.

\subsubsection{Proof of Proposition~\ref{p13:third:aligned-family}}
\label{p14:comp:third:aligned-family}

\begin{proof}
Fix one channel $a:b$ and orient the source pair $e=\{x,y\}$.
Choosing $B^{O_{a,b}(1)}$ offsets, write
$x=at-i_0$, $y=bt-j_0$ and $X=abt$.
Merge repeated values in $bV_x\cup aV_y$ to obtain
$U_X\subset X+[-D B,D B]$, where $D$ depends only on the fixed channel.
Use radius $DB$ for alignment with this compressed source, and choose
the large-prime cutoff above $(2D+1)B$.  The old centre lies in a fixed
expanded interval starting at $N$, while the new centre remains in $I_N$,
as required by the expanded version of Lemma~\ref{p1:arith:asymmetric}.
Multiplying an even selection by a constant does not change its square
class, so
\begin{equation}
 S_x+S_y\subseteq
       \operatorname{span}\{v(n):n\in U_X\}.
 \label{p13:third:true-compression}
\end{equation}
This inclusion holds regardless of injectivity on $U_X$.
The uniform aligned rank bound of
Lemma~\ref{p1:arith:alignedrank} gives
\begin{equation}
 q(e)\leq N^{-\kappa_b+o(1)},\qquad
                   \kappa_b=2-\frac1b\geq\frac53.
 \label{p13:third:source-weight}
\end{equation}
For fixed height the residual defect is uniform, so the same bound
includes admissible local rank defects.

\paragraph{Targets not aligned with the compressed source.}
For $z\notin e$, put
$\delta=\dim(S_z\cap(S_x+S_y))$.
Connectedness of $T=e\cup\{z\}$ implies $\delta>0$.
The inclusion \eqref{p13:third:true-compression} bounds $\delta$ by
the asymmetric rank loss $\delta^*(X,z)$ of
Lemma~\ref{p1:arith:asymmetric}.  That lemma permits an arbitrary true
set $U_X$ of this fixed logarithmic width; it does not require injectivity.
It gives
\[
 \sum_{X,z\ {\rm nonaligned}}(2^{\delta^*(X,z)}-1)
                         \ll N^{31/24+\varepsilon'}.
\]
Since $q(T)=q(e)p_z2^\delta$, $p_z\leq2p$, and
$2^\delta\leq2(2^{\delta^*}-1)$, these targets have total mass at most
\begin{equation}
 N^{-\kappa_b-1+31/24+\varepsilon'+o(1)}
                     \ll N^{-11/8+\varepsilon}.
 \label{p13:third:unaligned-target-cost}
\end{equation}
The number of source-offset patterns is polynomial in $B$ and is
absorbed into the exponent margin.

\paragraph{The residual kernel for an aligned target.}
This sector includes every target aligned at the same enlarged radius
$DB$.  Enlarging that fixed radius changes only the constants in the
polynomial offset and height bounds below.
Write an aligned target as
$z=(ct+h)/d$, with $(c,d)=1$ and
$c,d,|h|\leq B^{O_{a,b}(1)}$, as in the proof of
Lemma~\ref{p1:arith:fixed-channel-rank}.
Before this final normalization, $U_X$ and $V_z$ each have fixed
logarithmic width.  Choose $H>C_D B$ larger than both diameters and
remove the prime rows at most $H$ and those dividing the normalization
multipliers.  Each remaining prime row meets at most one distinct old
value and at most one new value.  This uses the fact that repeated
old values were merged before the target was added.

After the common normalization, merge equal values once more.
There are at most $3B$ distinct integers, all at height at least a
constant multiple of $N$ and in an interval of diameter
$B^{O_{a,b}(1)}$.  The remaining prime rows have weight at most two,
and leave one variable per unpinned component.  We give the uniform
residual estimate, since the coefficients $c,d$ need not be fixed.

For a fixed integer $K$, restrict to components of at most $K$ vertices.
The removed prime rows and a total-parity row form a code on these
components of redundancy
\[
 s\leq\pi(H)+O_{a,b}(\log B)+1=O_{a,b}(B/\log B).
\]
A nonzero word in its kernel gives an even square product of distinct
normalized integers.  Lemma~\ref{p1:arith:runge}, applied with the
polynomial-in-$B$ diameter just specified, forces it to contain
$\gg_{a,b}\log N/\log B$ physical vertices and hence
$\gg_{a,b,K}\log N/\log B$ components.
If $c_K$ is the number of these components, the Hamming inequality with
$t\asymp_{a,b,K}\log N/\log B$ gives
\[
 \sum_{j=0}^{t}\binom{c_K}{j}\leq2^s,
 \qquad c_K=O_{a,b,K}(B/\log B).
\]
Here $t$ is chosen below half the component minimum distance.
Components larger than $K$ number at most $3B/K$.
Choosing $K$ arbitrarily large before $N$ therefore bounds the entire
value kernel on the distinct normalized set by $o_{a,b}(L)$,
uniformly over $c,d,h$.

The relevant map factors as
\[
 \bigoplus_{v\in T}E_v
   \longrightarrow\{\text{merged normalized coordinates}\}
   \longrightarrow\{\text{square classes}\}.
\]
Its nullity is at most the nullity of merging plus the residual $o(L)$
just proved.  Internal defects of the old windows are included in this
composition and do not require any assumption of good normalization.

\paragraph{Counting the exact coincidences.}
In the common parameter the physical coefficients are $ad,bd,c$.
Suppose first that they are distinct.  Divide by their common gcd
and order them as $u<v<w$.  Unless the parameter belongs to a
$B^{O_{a,b}(1)}$ endpoint population, dyadic geometry gives $w<2u$.
Indeed a rational ratio of polynomial height cannot cross two by
$B^{O(1)}/N$, and equality forces the extreme windows to the two
dyadic endpoints.
The lattice calculation \eqref{p13:third:density} and
Lemma~\ref{p13:third:rational-rank} show that merging loses at most
$19L/30+O(1)$ dimensions.  The residual estimate therefore gives
\[
 \dim(S_x+S_y+S_z)\geq\frac{71}{30}L-o(L),\qquad
 q(T)\leq N^{-71/30+o(1)}.
\]
There are $NB^{O_{a,b}(1)}$ aligned triples of this form, including the
choices of offsets and rational coefficients.  Their total mass is
$O(N^{-41/30+\varepsilon})$.

If two physical coefficients are equal, compatibility makes the
interiors of the corresponding parameter intervals disjoint, since
distinct compatible starts have separation at least $L$.
Their combined coincidences with the third lattice are bounded by the
length of that third interval: at most $L/a+O(1)$ or $L/b+O(1)$ units.
As $a\geq2$, the total merging nullity is at most $L/2+O(1)$.
Including the residual, this gives mass $O(N^{-3/2+\varepsilon})$.
Endpoint populations have only $B^{O_{a,b}(1)}$ source and target
choices; the bound $q(T)\leq q(e)$ from
\eqref{p13:third:source-weight} gives $O(N^{-5/3+\varepsilon})$.
These estimates include exact contacts.  Strict contacts are incompatible
and do not enter the sum.

The least of the exponents obtained is $41/30$.  Sum over the fixed
finite family $\mathcal F$ to conclude.
\end{proof}


\paragraph{Homogeneous variant for separated triples.}
The same $O_{C,\mathcal F,\varepsilon}(N^{-41/30+\varepsilon})$ bound
holds for the sum of
$q_0(T)=2^{-\dim\sum_{x\in T}S_x}$ over connected triples
$T\subset G_N$ with distinct, pairwise separated starts and containing
a source pair in $\mathcal A_{\rm raw}$ in one of the fixed channels.
Here $q_0$ is the homogeneous probability, even when the original affine
prescriptions are incompatible.

To see this, repeat the proof above with $q_0$ in place of $q$.
The source estimate \eqref{p13:third:source-weight} is a rank estimate
and holds for $q_0(e)$; the target marginal is exactly $p$ on $G_N$.
The nonaligned identity is
$q_0(T)=q_0(e)p2^\delta$, giving the same $N^{-11/8+\varepsilon}$ bound.
In the aligned sector all residual-kernel and merging estimates are
phase-free. When two physical coefficients coincide, the imposed
separation makes the corresponding interval interiors disjoint, so
the $N^{-3/2+\varepsilon}$ estimate still holds without compatibility.
The distinct-coefficient estimate is $N^{-41/30+\varepsilon}$, and the
endpoint estimate follows from $q_0(T)\le q_0(e)$ as before.
Exact contacts at distance $L$ remain allowed. This proves the stated
homogeneous variant independently of the moment and reconstruction theorems.

\subsection{Positive Palm costs on all target windows}

A good edge changes the law of a third start by introducing an additional
intersection of square-class spaces. Its positive effect can be bounded
by the same positive triple mass, even when the target is locally irregular.

\begin{proposition}[Positive Palm cost of a fixed channel]
\label{p1:arith:palm}
\label{p1:arith:palm-all-targets}
Fix a reduced interior channel $a:b$, with $0<a<b<2a$, and an integer
$D\ge6b$. Choose an integer private-prime threshold with
$6bB\le H\le DB$.
Let $\mathcal E_{a,b}$ be its compatible dependent good edges, write $J_e=\prod_{x\in e}J_x$ and
$q_e=\P(J_e=1)$, and put $\beta_{a,b}=\sum_{e\in\mathcal E_{a,b}}q_e$.
For a deterministic set of targets $G\subseteq I_N$, define
\begin{align*}
 M_{\mathrm{ex}}(G)
   &=\sum_{e\in\mathcal E_{a,b}}\sum_{z\in G\setminus e}
       \bigl(q(e\cup\{z\})-q_ep_z\bigr)_+,\\
 M_{\mathrm{nom}}(G)
   &=\sum_{e\in\mathcal E_{a,b}}q_e\sum_{z\in G\setminus e}
       \bigl(\P(J_z=1\mid J_e=1)-p\bigr)_+.
\end{align*}
For every fixed critical band and $\varepsilon>0$, uniformly in $G$
and in the stated range of $H$, for all sufficiently large $N$,
\begin{equation}
 \begin{split}
 &\beta_{a,b}>0,\qquad
 M_{\mathrm{ex}}(G)+M_{\mathrm{nom}}(G)
     \ll_{a,b,C,D,\varepsilon}N^{-41/30+\varepsilon},\\
 &\frac{M_{\mathrm{ex}}(G)+M_{\mathrm{nom}}(G)}{\beta_{a,b}}
     \ll_{a,b,C,D,\varepsilon}N^{-11/30-1/b+\varepsilon}.
 \end{split}
 \label{p2:palm:all-targets-bound}
\end{equation}
No regularity condition is imposed on the target windows.
\end{proposition}

\begin{proof}
For $z\notin e$, put
$d_e(z)=\dim(S_z\cap\sum_{x\in e}S_x)$.
The finite affine rank bound gives
\[
 q(e\cup\{z\})\le q_ep_z\,2^{d_e(z)}.
\]
The inequality includes incompatible joint prescriptions and zero
marginals. Thus $p_z=0$ forces $q(e\cup\{z\})=0$, and $d_e(z)=0$
makes the exact positive excess vanish.
A positive summand of $M_{\mathrm{ex}}$ therefore has
$d_e(z)>0$ and $q(e\cup\{z\})>0$.
Compatibility forces all three pairwise distances to be at least $L$,
including possible equality. The source is a dependent pair, so neither
source singleton can split off as a direct component. Positive crossing
prevents the target singleton from splitting off. The triple is connected.
Each summand is at most $q(e\cup\{z\})$, and a triple has at most three
choices of a marked source pair. Proposition~\ref{p13:third:aligned-family},
which includes irregular windows and exact contacts, gives
\[
 M_{\mathrm{ex}}(G)
 \le3\sum_{\substack{|T|=3,\ T\ \mathrm{connected},\ q(T)>0\\
              T\text{ contains a pair in the fixed channel}}}q(T)
 \ll N^{-41/30+\eta}
\]
for every fixed $\eta>0$.

All conditioned sources have $q_e>0$. The positive-part inequality gives
\[
 \bigl(q(e\cup\{z\})-q_ep\bigr)_+
 \le\bigl(q(e\cup\{z\})-q_ep_z\bigr)_+
      +q_e(p_z-p)_+.
\]
Using the marginal estimate~\eqref{p1:arith:marginalsum},
\[
 M_{\mathrm{nom}}(G)
 \le M_{\mathrm{ex}}(G)+\beta_{a,b}\sum_{z\in I_N}(p_z-p)_+
 \ll N^{-41/30+\eta}+\beta_{a,b}N^{-1+\eta}.
\]
Lemma~\ref{p1:arith:fixed-channel-mass} gives, in the stated cutoff range,
$\beta_{a,b}\asymp_{a,b,C,D}LN^{-1+1/b}>0$.
Each unordered compatible edge has two ordered orientations, so its
mass is half the ordered mass in that lemma. Monotonicity of the
private-prime condition sandwiches the mass between the phases at
$H=6bB$ and $H=DB$, proving uniformity throughout the stated range.
Since $b\ge3$, the second raw term is
$O(LN^{-2+1/b+\eta})$, smaller than the $41/30$ scale.
Division by the positive phase bounds the normalized sum by a constant
times
\[
 L^{-1}N^{-11/30-1/b+\eta}+N^{-1+\eta}.
\]
The second exponent is smaller because $11/30+1/b\le7/10<1$.
Choose the internal margin within $\varepsilon$ and enlarge the scale
threshold to absorb fixed factors. All parameters precede $N,L,H,G$.
\end{proof}

For the two principal channels, one common choice $24B\le H\le DB$
with fixed $D\ge24$ gives
\begin{align}
 \frac{M_{\mathrm{nom},2:3}(G)}{\beta_{2,3}}
     &\ll_{C,D,\varepsilon}N^{-7/10+\varepsilon},
     \label{p1:arith:palm23}\\
 \frac{M_{\mathrm{nom},3:4}(G)}{\beta_{3,4}}
     &\ll_{C,D,\varepsilon}N^{-37/60+\varepsilon}.
     \label{p1:arith:palm34}
\end{align}
These are bounds for positive increases in the target probabilities;
no optimality or absolute-discrepancy statement is asserted.
